% ID: FIS-ELE-COU-001
% Titolo: Bilancia di torsione di Coulomb
% Disciplina: Fisica
% Area: Elettrostatica
% Argomento: Legge di Coulomb
% Sottoargomento: Verifica sperimentale della forza elettrica
% Classe: Quinta liceo scientifico
% Difficolta: 2
% Tipo: quesito_teorico
% Risultato: verifica la dipendenza F proporzionale a q_1 q_2 / r^2
% Tempo_stimato: 6 min
% Competenze: descrizione di un esperimento, collegamento teoria-esperimento, argomentazione
% Prerequisiti: forza elettrica, carica elettrica, equilibrio
% Tag: fisica, elettrostatica, legge_coulomb, bilancia_torsione, forza_elettrica
% Fonte: verifica_fisica_elettrostatica_anonimizzata
% Autore: Riccardo Carli
% Licenza: CC-BY-SA-4.0

\begin{esercizio}
Descrivi l'esperimento della bilancia di torsione di Coulomb e spiega quale
relazione fisica permette di verificare.
\end{esercizio}

\begin{soluzione}
La bilancia di torsione di Coulomb e costituita da un'asta leggera sospesa a un
filo sottile. Su un'estremita dell'asta si trova una piccola sfera carica. Una
seconda sfera carica viene avvicinata alla prima.

La forza elettrica tra le due sfere fa ruotare l'asta e torce il filo. Il filo
reagisce con una coppia elastica che aumenta con l'angolo di torsione. Quando
il sistema e in equilibrio, la coppia elastica del filo bilancia la coppia
dovuta alla forza elettrica.

Misurando l'angolo di torsione si puo risalire al modulo della forza elettrica.
Ripetendo l'esperimento a distanze diverse si verifica che la forza diminuisce
con il quadrato della distanza:
\[
F\propto \frac{1}{r^2}.
\]
Variando le cariche si trova inoltre che la forza e proporzionale al prodotto
delle cariche:
\[
F\propto q_1q_2.
\]

La relazione verificata e quindi la legge di Coulomb:
\[
F=k\frac{|q_1q_2|}{r^2}.
\]
\end{soluzione}
